\section{Entropy localization and the support profile}
\label{prof:entropy}

We now convert the posterior estimates into an inequality for every test
function of small support.  All final typical-environment assertions concern
the original, unplanted disorder.  The passage from the planted estimates to
these statements will be made before choosing a test function; this order is
essential for the uniformity of the resulting profile.

Write $\mu=\mu_J$ for the critical zero-field Gibbs measure.  For a probability
measure $\pi$ on the cube, use the unscaled heat-bath form
\[
 \mathcal D_\pi(f)
 =\sum_{i=1}^n\pi\!\left[\operatorname{Var}_\pi
                  (f\mid X_{-i})\right].
\]
Thus $\mathcal D_J$ is the form at $\mu_J$, and $\mathcal D_b$ will denote the
same form for the posterior law at observation parameter $b$.
When used for heat-bath dynamics, this unscaled form corresponds
to rate one at each site. The observation parameter \(b\) instead
measures Gaussian information with the precision rate defined below.

We use the following static conclusions established above, with exactly their
stated probability qualifications: the sharp small-time covariance estimate
\eqref{prof:eq22}, the estimate below the logarithmic threshold
\eqref{prof:eq25}, and the positive-time assertions in
Section~\ref{prof:positive-inputs}.  The latter include bounded covariance on
each compact positive-time interval, a terminal observation gap, the
stable-field implication \eqref{prof:eq26}, and the square-reweighting
estimates \eqref{prof:eq27} and \eqref{prof:eq27-coarse}.

In particular, the small-positive-time failure
rate for stability can be any prescribed fixed multiple of $n b^{3/2}$, and
the probability of an approximate zero with $q>u$ is at most
$\exp(-cnu^3)$ for \(C\sqrt b\le u\le u_*\), with a fixed
small \(u_*>0\). Constants and residual tolerances may
depend on a fixed compact positive-time interval.  The proofs in
Section~\ref{prof:positive-inputs} supply these statements at the critical
parameter itself; no subcritical spectral-gap theorem is being applied at
its endpoint.

\subsection{The observation identities}

The use of Gaussian observations here is a stochastic localization argument
in the sense of Eldan~\cite{Eldan2013StochasticLocalization}.
The variance and entropy transfers are related to the localization methods
of Eldan, Koehler, and Zeitouni~\cite{EldanKoehlerZeitouni2022} and the
framework of Chen and Eldan~\cite{ChenEldan2025Localization}.
We give the finite-state identities and the estimates needed here explicitly.

Choose the sufficiently large fixed terminal time $T$ supplied by the static
inputs.  On the spectral event $\|J\|\le3$, fix $c_+\ge3$ and put
$R=J+c_+I\succeq0$.  Continue the scalar observations after $T$ for one unit
of time with precision $R$.  Set $E=T+1$ and
\[
 \Lambda_b=I\quad(0\le b\le T),\qquad
 \Lambda_b=R\quad(T<b\le E).
\]
More explicitly, conditional on $X$, the field has independent increments
with drift $\Lambda_bX\,db$ and covariance $\Lambda_b\,db$.  The accumulated
quadratic term in its likelihood removes the interaction by time $E$:
at $b=T+\lambda$, the off-diagonal interaction is $(1-\lambda)J^\circ$, and
$\mu_E$ is a product law.

For a real function $f$ with $\mu(f^2)=1$, let $\nu=f^2\mu$.  Use $\mu_b$
and $\nu_b$ for the conditional spin laws given the observation history,
under the respective initial laws. For a fixed $J$, write
$\mathbb P_J^\pi$ for the observation-history law with initial spin law
$\pi$, and $\mathbb E^\pi$ for expectation under this conditional law.
Define
\[
 \begin{gathered}
 w_b=\mu_b(f^2),\qquad m_b^{\mathrm{post}}=\mu_b(X),\qquad
 \Sigma_b=\operatorname{Cov}_{\mu_b}(X),\\
 a_b=\nu_b(X)-\mu_b(X),\qquad
 H_b=\mathbb E^\nu\operatorname{KL}(\nu_b\Vert\mu_b).
 \end{gathered}
\]
Histories with $w_b=0$ have zero probability under $\nu$ and can be omitted.
The likelihood ratio of the two field-history laws is $w_b$, and
$\nu_b=f^2\mu_b/w_b$.  Consequently
\[
 H_b=\mathbb E^\mu\operatorname{Ent}_{\mu_b}(f^2),\qquad
 \operatorname{KL}(\text{field law under }\nu
                   \Vert\text{field law under }\mu)=H_0-H_b.
\]

The Gaussian filtering identities give, at every time away from the
precision change and in integral form across it,
\begin{equation}
 \begin{split}
 H'_b&=-\frac12\mathbb E^\nu
                  \|\Lambda_b^{1/2}a_b\|^2,\\
 a_b&=\mathbb E^\nu\left[
          a_E+\int_b^E\Sigma_v\Lambda_v a_v\,dv
          \,\middle|\,\mathcal F_b\right].
 \end{split}
 \label{prof:eq28}
\end{equation}
Here $\mathcal F_b$ is the observation history, including $J$.  To verify
the first identity, the field drifts under the two laws differ by
$\Lambda_ba_b$; the expected logarithm of their likelihood ratio therefore
has derivative $\mathbb E^\nu\|\Lambda_b^{1/2}a_b\|^2/2$.

Subtract this from the preceding entropy decomposition.  For the second
identity, the $\nu$-posterior mean is a martingale under $\nu$, whereas the
$\mu$-posterior mean has drift $\Sigma_b\Lambda_ba_b$ under that law.
Integrating their difference to $E$ proves the assertion.  These arguments
also follow directly by It\^o's formula for the finite posterior sums.

Two further identities will be used repeatedly.  Since $\mu_E$ is product,
binary Pinsker on each coordinate and entropy tensorization from below give
\begin{equation}
 \mathbb E^\nu\|a_E\|^2\le2H_E.
 \label{prof:product-mean}
\end{equation}
Also, conditioning additionally on the observations can only decrease an
averaged conditional variance.  Applied separately with $X_{-i}$ fixed, this
gives
\begin{equation}
 \mathbb E^\mu\mathcal D_b(f)\le\mathcal D_J(f),\qquad
 \mathbb E^\nu\frac{\mathcal D_b(f)}{w_b}
       =\mathbb E^\mu\mathcal D_b(f)\le\mathcal D_J(f).
 \label{prof:energy-under-observation}
\end{equation}

\subsection{A profile with a bounded test likelihood}

\begin{proposition}[Restricted support profile]
\label{prof:restricted-profile}
For every fixed $\omega>0$ there are $c_\omega>0$ and disorder events
$\mathcal G_n(\omega)$ with probability tending to one such that, on
$\mathcal G_n(\omega)$, simultaneously for every real $f$ satisfying
\[
 \mu(f^2)=1,\qquad
 0<p:=\mu(f\ne0)\le\tfrac12,\qquad
 \|f\|_\infty^2\le100p^{-2},
\]
one has
\begin{equation}
 \mathcal D_J(f)\ge c_\omega n^{-2/3-\omega}\log(1/p).
 \label{prof:eq29}
\end{equation}
\end{proposition}

\begin{proof}
Choose the covariance accuracy \(\epsilon>0\) sufficiently small
in terms of the requested \(\omega\), and obtain the corresponding
\(a,s_0,K_{\log}\) from \Cref{prof:prop-amplified-covariance}.
This same accuracy is used in all three entropy ranges below.
Put $h=\log(1/p)$.  Entropy relative to the support and the upper bound on
$f^2$ imply
\begin{equation}
 h\le H_0=\mu(f^2\log f^2)\le2h+\log100.
 \label{prof:initial-entropy}
\end{equation}
On $\|J\|\le3$, every atom of $\mu$ is at least
$2^{-n}e^{-3n}$, so $h\le(3+\log2)n$.  We treat three ranges of $h$.

\paragraph{Moderate entropy.}
For constants \(K'\) and \(c'\) fixed in the order below, first suppose
\[
 K'\log n\le h\le c'n,
 \qquad r_h=(h/n)^{1/3},\qquad t_h=r_h^2.
\]
The environment restrictions below are imposed on dyadic bins for $h$;
using an endpoint of a bin in place of $h$ only changes fixed constants.
Fix the low-likelihood exponent \(C_2>0\) and a desired negligible
error order \(n^{-M}\). Choose \(C_{\rm rate}\) with fixed slack
beyond the likelihood exponent \(4+C_2\) and the subsequent
Markov losses. The remaining constants will be chosen after the
following static consequences are stated.
For every required fixed rate $C_{\mathrm{rate}}$, the preceding static
estimates have the following consequences. All integrated probabilities
in this list use planted disorder and reference observations with
the spectral restriction imposed; the first also draws two
conditionally independent replicas from \(\mu_b\).
\begin{enumerate}
\item For $0\le b\le c_1t_h$, where $c_1\le1$, a reference posterior pair
has $|\langle X^1,X^2\rangle|\le C_1nr_h$, except with integrated
probability at most $e^{-C_{\mathrm{rate}}h}$.

Here $C_1$ need not increase
as $c_1$ decreases.
\item On each fixed-ratio window $[c_1t_h,L_1t_h]$, with
$0<c_1<L_1<\infty$, the replica log-moment estimate
\eqref{prof:eq5} holds with time $b$ and scale $nb^{3/2}$, with the required
small off-diagonal tilts and any prescribed fixed relative error, outside
integrated probability $e^{-C_{\mathrm{rate}}h}$.
\item By choosing $L_1$ large enough, the sharp bound
$\|\Sigma_b\|\le(1+\epsilon)/b$ on $[L_1t_h,s_0]$ has the same failure
rate.  Beyond $s_0$, use the bounded positive-time covariance, including the
terminal extension.  At $E$, the product field satisfies
$\max_i|(h_E)_i|\le C'\sqrt h$ with that rate, also when the exceptional
probability is weighted by $1+\max_i|(h_E)_i|$.
\end{enumerate}

For completeness, the first assertion does not need a lower spherical
width bound at very early times.  Impose the upper width bound at scale
$t_h$ with an arbitrarily large fixed rate.  The spherical length-deficit
estimate then controls the pair integral relative to $Z_o(h_b)^2$, using
\eqref{prof:eq6}, the field-norm bound, and the large-overlap comparison.
Exact shell conditioning costs only a polynomial by the spectral and
spherical estimates; translations improve the required small-ball bound.

The trace loss is estimated at shifts above the upper width bound.
Orientations with $Z(h_b)/Z_o(h_b)<e^{-C_{\mathrm{rate}}h}$ have at most
that planted mass conditional on the eigen-coordinates, so division by the
partition function is absorbed by an increase of the rate budget.  The
second and third assertions follow from the spherical width tails,
\eqref{prof:eq5}, \eqref{prof:eq22}, and the positive-time inputs.  For the
last field maximum, Gaussian tails and a union bound apply; under planting,
the $RX$ term has Gaussian site marginals with bounded translates.  Gaussian
moment bounds give the assertion with the extra maximum-field factor.

Complete the fixed parameter choices before constructing the disorder
event. Obtain \(C_1\) at the chosen rate and then take \(c_1\) small enough
for \eqref{prof:early-entropy}; \(C_1\) does not increase under this choice.
Next enlarge \(L_1\). At the start of its sharp window the covariance
exponent is
\[
 an(L_1t_h)^{3/2}=aL_1^{3/2}h,
\]
so \(L_1\) can pay the required rate. Then choose the fixed replica number
and its relative accuracy for the intermediate window. Finally decrease
\(c'\): a positive-time exception \(e^{-cn}\) has exponent at least
\((c/c')h\), and this choice also keeps \(L_1t_h\le s_0\) and all scales
in their allowed ranges. Increase \(K'\) to absorb polynomials and the
logarithmic threshold. We may require all discarded errors below to be
$O(n^{-M})$ for the fixed sufficiently large $M$ chosen above.

With these choices fixed, choose the disorder event on which the support
profile will hold. Entropy levels have been placed in dyadic bins; in each
bin integrate the conditional exceptions over its three time windows.
This construction uses the bin and fixed tolerances, before any test
function is chosen. For a field failure \(B_b\), write
\(q_J(b)=\mathbb P_J^\mu(B_b)\) for its reference-field probability
conditional on \(J\). For a pair failure, \(q_J(b)\) denotes the
conditional field average of its \(\mu_b^{\otimes2}\) probability.
The ordinary high-probability
lower bound $Z(0)\ge n^{-C}e^{n/4}$, together with
$\mathbb EZ(0)=e^{n/4}$, transfers planted exceptions on the spectral event
to ordinary exceptions at polynomial cost. Indeed, with
\(L(J)=Z(0)/\mathbb EZ(0)\),
\[
 \E_{\rm ord}\!\left[\mathbf1_{\{\mathcal G,\,L\ge n^{-C}\}}
                    q_J(b)\right]
 \le n^C\E_{\rm pl}[\mathbf1_{\mathcal G}q_J(b)].
\]
Integrate over each specified time window, apply Markov with unused rate
slack, and take a union over the $O(\log n)$ bins.

This produces one typical disorder event, independent of
$f$, on which all required conditional time-integral exceptions hold.
There is enough unused rate to absorb the later likelihood factors.
The chosen $K'$ makes their products with any specified polynomial in $n$
and $e^{C_0h}$ negligible, for each of the finitely many constants $C_0$
used below. Endpoint tests are treated in the same way.

On the early window, the pair bound implies
\begin{equation}
 \mathbb E^\nu\|a_b\|^2\le Cnr_h
                  +\text{an error negligible in time integral}.
 \label{prof:early-drift}
\end{equation}
Indeed the square of a posterior mean is bounded by the posterior
expectation of the absolute pair inner product.  To apply this also to
$\nu_b^{\otimes2}$, discard $w_b<e^{-C_2h}$: its probability under $\nu$
is at most $e^{-C_2h}$, and its squared-drift contribution is at
most \(4ne^{-C_2h}\). On the complement, the density of the
\(\nu\)-field and two \(\nu_b\) replicas relative to the
reference field and two \(\mu_b\) replicas is exactly
\[
 w_b\,\frac{f^2(X^1)}{w_b}\frac{f^2(X^2)}{w_b}
 =\frac{f^2(X^1)f^2(X^2)}{w_b}
 \le10^4e^{(4+C_2)h}.
\]

The ordinary posterior mean evaluated under the
$\nu$-field law is controlled in the same way.  Taking $c_1>0$ sufficiently
small, \eqref{prof:eq28}, \eqref{prof:initial-entropy}, and
$nr_ht_h=h$ now give
\begin{equation}
 H_{c_1t_h}\ge h/2.
 \label{prof:early-entropy}
\end{equation}

On the intermediate window, use $l$ replicas in \eqref{prof:eq5} with
off-diagonal matrix entries $\delta/l$, where $\delta>0$ is sufficiently
small and independent of $l$.  The entropy variational inequality applied
to $\nu_b^{\otimes l}$ gives a multiple of
$b\|\nu_b(X)-m_o\|^2$ bounded by a constant times
$\operatorname{KL}(\nu_b\Vert\mu_b)$ and the log-moment cost divided by
$\delta(l-1)$.  Indeed the expectation of the tilt is exactly
$b\delta(l-1)\|\nu_b(X)-m_o\|^2/2$, whereas the relative entropy is
$l\operatorname{KL}(\nu_b\Vert\mu_b)$.

The squared Hilbert--Schmidt norm
of the tilt matrix is bounded by $\delta^2$, uniformly in $l$.  Apply the
same inequality to ordinary replicas, whose entropy cost is zero, and use
$\|a_b\|^2\le2\|\nu_b(X)-m_o\|^2+2\|\mu_b(X)-m_o\|^2$.
After averaging over the $\nu$-field law, this yields
\begin{equation}
 b\mathbb E^\nu\|a_b\|^2
       \le C_\delta H_b+\zeta nb^{3/2}
                  +\text{negligible integrated error}.
 \label{prof:intermediate-drift}
\end{equation}
Here $\zeta>0$ can be as small as desired by first taking a sufficiently
large fixed $l$ and then the relative log-moment error sufficiently small.

All exceptional field likelihoods are paid for by the same sup-norm bound.
Gronwall on this fixed-ratio window, followed by a sufficiently small
choice of $\zeta$, gives
\begin{equation}
 H_{L_1t_h}\ge c h.
 \label{prof:middle-entropy}
\end{equation}
The positive constant may depend on the fixed ratio $L_1/c_1$, which causes
no loss of a power of $n$.

For later times let $c_v$ be the good covariance bound:
$c_v=(1+\epsilon)/v$ up to $s_0$, and a fixed constant afterwards.  Set
$M_v=\mathbb E^\nu\|\Lambda_v^{1/2}a_v\|^2$.  Conditional Jensen in
\eqref{prof:eq28}, then Cauchy--Schwarz including its endpoint term, gives
\begin{equation}
 \begin{split}
 (\mathbb E^\nu\|a_b\|^2)^{1/2}
 &\le \sqrt{2H_E}
       +\int_b^E c_v\|\Lambda_v\|^{1/2}\sqrt{M_v}\,dv+n^{-M}\\
 &\le \sqrt{2H_b}
       \left(1+\int_b^E c_v^2\|\Lambda_v\|\,dv\right)^{1/2}
       +n^{-M}.
 \end{split}
 \label{prof:backward-entropy}
\end{equation}
We used $2H_E+\int_b^E M_v\,dv=2H_b$.

The deterministic bounds
$\|\Sigma_v\|\le n$, $\|a_v\|\le2\sqrt n$ control the discarded
integrals at polynomial cost.  Squaring introduces only another negligible
error, since all remaining factors are polynomially bounded on this range.
As $\Lambda_b=I$ below $s_0$, the coefficient in the resulting differential
inequality for $H_b$ is at most
\[
 \frac{(1+\epsilon)^2}{b}+C\qquad(L_1t_h\le b\le s_0),
\]
and is bounded afterwards.  Gronwall and \eqref{prof:middle-entropy} imply,
with $\kappa=(1+\epsilon)^2$,
\begin{equation}
 H_E\ge c t_h^\kappa h-n^{-M+O(1)}.
 \label{prof:entropy-at-product}
\end{equation}

For a Bernoulli law with smaller atom $z$, the elementary two-point
inequality is
\[
 \operatorname{Ent}(u^2)
       \le C\log(1/z)\operatorname{Var}(u).
\]
One may first replace $u$ by $|u|$.  For comparable nonnegative values use
$\operatorname{Ent}(u^2)\le\operatorname{Var}(u^2)/\mathbb E u^2$;
if the values are not comparable, bound the entropy by the larger value
squared times binary entropy and compare the difference of the values.
For a spin of field $v$, $\log(1/z)\le C(1+|v|)$.  Tensorizing at the
product endpoint and using \eqref{prof:energy-under-observation} therefore
gives
\begin{equation}
 H_E\le C\mathbb E^\mu\left[
       (1+\max_i|(h_E)_i|)\mathcal D_E(f)\right]
       \le C'\sqrt h\,\mathcal D_J(f)+n^{-M}.
 \label{prof:product-entropy}
\end{equation}
To justify the correlated exceptional term, put
\(B_E=\{\max_i|(h_E)_i|>C'\sqrt h\}\). Pointwise,
\(\mathcal D_E(f)\le n\mu_E(f^2)=nw_E\le n\|f\|_\infty^2\).
Hence its contribution is at most
\begin{equation}\label{prof:product-exception}
 n\|f\|_\infty^2\,
 \E^\mu\!\left[(1+\max_i|(h_E)_i|)\mathbf1_{B_E}\right].
\end{equation}
The weighted endpoint exception and the restricted supremum bound
make this negligible. On \(B_E^c\), pull out \(C'\sqrt h\) and
use \eqref{prof:energy-under-observation}.

Combining \eqref{prof:entropy-at-product} and
\eqref{prof:product-entropy} yields
\[
 \mathcal D_J(f)\ge c n^{-2\kappa/3}h^{1/2+2\kappa/3}.
\]
At $\kappa=1$ the right side is $c n^{-2/3}h^{7/6}$.
The initial choice of \(\epsilon\) in terms of \(\omega\) proves
\eqref{prof:eq29} in this range, since $h\ge1$ and the extra exponent of
$n$ can be less than $\omega$.

\paragraph{Small entropy.}

In this range the test likelihood is only polynomially large, so the
bottom-scale estimate can absorb its exceptional contribution. We use
variance rather than entropy; the accumulated covariance bound has the
leading exponent needed for the support profile.
For $h<K'\log n$, use
$V_b=\mathbb E^\mu\operatorname{Var}_{\mu_b}(f)$ up to $T$.
Cauchy--Schwarz on the support gives $V_0\ge1-p\ge1/2$.
Filtering and the covariance bound imply
\[
 -V'_b=\mathbb E^\mu
          \|\operatorname{Cov}_{\mu_b}(f,X)\|^2
       \le c_bV_b+\text{negligible integrated error}.
\]
Here the exceptional errors can be made arbitrarily small polynomially:
the test bound is polynomial in $n$, \eqref{prof:eq25} has any required
inverse-polynomial failure rate, and \eqref{prof:eq22} is used only above
a sufficiently large logarithmic threshold.

Specifically take
$b_*=(K''\log n/n)^{2/3}$, with $K''$ enlarged as necessary.  Below $b_*$
use $c_b\le Cn^{2/3}(\log\log(n+1000))^C$ from
\eqref{prof:eq25}; its integral is $o(\log n)$.  Between $b_*$ and $s_0$
use $(1+\epsilon)/b$, and then the bounded positive-time input.  Thus
\[
 \int_0^T c_b\,db
     \le\tfrac23(1+\epsilon)\log n+o(\log n).
\]
The same transfer and time-integrated Markov argument supplies one typical
environment event for these assertions.  Gronwall gives
$V_T\ge n^{-2/3-O(\epsilon)-o(1)}$.  The terminal gap and
\eqref{prof:energy-under-observation} give
$V_T\le C\mathcal D_J(f)+n^{-M}$. The initial choice of $\epsilon$ leaves
enough power slack to absorb $h\le K'\log n$ and obtain
\eqref{prof:eq29}.

\paragraph{Entropy of order $n$.}
It remains to treat $h>c'n$. The maximum test likelihood no longer
controls the fixed positive-time exceptions. Instead, we charge those
exceptions to the relative entropy spent in changing the observation law.
The resulting differential inequality forces that spent entropy to
vanish at leading order. We show that, for
every fixed $C_*>0$, with probability tending to one no test function in
this range can have $\mathcal D_J(f)\le C_*n^{1/3}$.
Taking $C_*=1$ implies the desired estimate because
$h\le(3+\log2)n$.

The active class throughout this paragraph is still
\[
 \mu(f^2)=1,\qquad p=\mu(f\ne0)\le\tfrac12,\qquad
 \|f\|_\infty^2\le100p^{-2},\qquad h=\log(1/p)>c'n.
\]
We choose the diagnostics before selecting a violating function.
Take a countable list of compact positive-time bands, integer rate
multipliers \(A_*\), the fixed dyadic outer-threshold grids from
\Cref{prof:positive-inputs}, and residual and mean tolerances
tending to zero on each band. For each fixed finite collection,
the integrated planted bounds, the size-bias lower bound, and
Markov with unused rate slack give an ordinary disorder event of
probability tending to one. Intersect it with the finite ordinary word
diagnostics, the original spectral restriction used to define \(R\), and
the ordinary disorder event in \Cref{prof:terminal-extension}. This
intersection still has ordinary probability tending to one, and all its
restrictions are independent of \(f\). On this event the corresponding
conditional reference-field bounds hold outside time sets whose
Lebesgue measure tends to zero.

If the desired conclusion failed with probability bounded away
from zero, diagonal selection would therefore give sizes tending
to infinity, environments satisfying successively more of these
finite collections, and a violating restricted function on each.
All diagnostics in this selection precede that function. Work
along such a sequence and define
\[
 e_n(b)=\frac{H_0-H_b}{n},\qquad 0\le b\le T.
\]
These functions vanish at zero, are nondecreasing, and are $2$-Lipschitz by
\eqref{prof:eq28} and $\|a_b\|\le2\sqrt n$.  Passing to a subsequence,
assume $e_n\to e$ uniformly.  We prove $e=0$.

For the fixed environment \(J\) and \(b\in[0,T]\), let
\(B\in\mathcal F_b\) be an event of the observation history through \(b\)
with \(\mathbb P_J^\mu(B)\le e^{-nw}\), where \(w>0\). The relative
entropy of the restrictions of the two field laws to \(\mathcal F_b\),
conditional on this same \(J\), followed by binary relative entropy, gives
\begin{equation}
 \mathbb P_J^\nu(B)
   \le\frac{e_n(b)+n^{-1}\log2}{w}.
 \label{prof:binary-transfer}
\end{equation}
Indeed this restricted field relative entropy is $ne_n(b)$, and its projection to the
indicator of $B$ is at least $\mathbb P_J^\nu(B)nw-\log2$.

At small fixed positive times apply \eqref{prof:eq26} with any fixed large
rate multiplier $A_*$.  Its bad event has $w=A_*b^{3/2}$ in
\eqref{prof:binary-transfer}, after using rate slack in the environment
transfer. On its complement, \eqref{prof:eq26} implies the
whole-region stability at \(A=V_y\) required by
\eqref{prof:eq27}; that mean estimate uses no overlap lower bound.
It makes
$\mathbb E^\nu[\|a_b\|^2/n]$ arbitrarily small as $n\to\infty$ and
then its tolerance tends to zero.  To see the error scale explicitly, the
additive energy error in that estimate has expectation at most
\[
 C_\zeta n^{-1/2}(1+\mathcal D_J(f))=O_{\zeta,C_*}(n^{-1/6})
\]
by \eqref{prof:energy-under-observation}.  Since
$\|a_b\|/\sqrt n\le2$, this error is charged linearly when estimating
the squared drift.

On the bad event use the approximate-zero estimate
\eqref{prof:eq27-coarse}.
Write $Q_b$ for the supremum of $q(y)$ over the sufficiently-small-residual
points occurring there.  It yields, in squared units, an upper bound
$CQ_b$ plus the same vanishing error. For
\(C\sqrt b\le u\le u_*\), the outer estimate supplies
\(\mathbb P_J^\mu(Q_b>u)\le e^{-cnu^3}\).
On a compact positive-time band, and with finitely many thresholds, one
may use a single residual tolerance smaller than all those provided:
an approximate zero at this tolerance remains one at each larger tolerance.

Put \(R_b=C'A_*^{1/3}\sqrt b\), taking \(C'\) large enough
for the outer range and the initial interval short enough that
\(R_b\le u_*\). If \(B_b\) is the stability failure, then
\(0\le Q_b\le1\) gives
\[
\begin{split}
 \E^\nu[Q_b\mathbf1_{B_b}]
 \le{}&R_b\mathbb P_J^\nu(B_b)
 +\sum_{\substack{j\ge0\\2^jR_b<u_*}}
       2^{j+1}R_b\,\mathbb P_J^\nu(Q_b>2^jR_b)\\
 &+\mathbb P_J^\nu(Q_b>u_*).
\end{split}
\]
On a fixed compact band the moving thresholds are bracketed by
the finite fixed dyadic grid described in
\Cref{prof:positive-inputs}, losing at most a fixed factor in
their cubic rates. Binary transfer makes the first term at most
\(CR_be(b)/(A_*b^{3/2})\) in the limit, while the sum is at most
\[
 Ce(b)\sum_{j\ge0}
       \frac{2^jR_b}{(2^jR_b)^3}
 \le Ce(b)/R_b^2.
\]
The final saturated tail uses the rate at \(u_*\) and contributes
at most \(Ce(b)/u_*^3\le C_{u_*}e(b)/R_b^2\).
Thus no outer rate beyond the supported small range is used.
Consequently, in integral form on every compact band in the initial
interval,
\begin{equation}
 e'(b)\le C A_*^{-2/3}\frac{e(b)}{b}.
 \label{prof:macroscopic-drift}
\end{equation}
The constant \(C\) is independent of \(A_*\). To see this,
first fix \(A_*\), then take \(n\to\infty\), and then send the
mean tolerance to zero. The stability margins and the refined
mean constants, which may depend on \(A_*\), disappear in this
order. The remaining coefficient comes from the absolute leading
constant \(2\) in \eqref{prof:eq27-coarse}, the absolute outer-tail
constants, and the fixed \(u_*\).

Choose $A_*$ so large that
$\theta=CA_*^{-2/3}<1$.  For $0<a<b$ in this interval, Gronwall and
$e(a)\le2a$ imply
$e(b)\le2b^\theta a^{1-\theta}$; letting $a\downarrow0$ proves
$e(b)=0$.  On the remaining positive times, the fixed-rate version of
\eqref{prof:eq26}, the good-field estimate \eqref{prof:eq27}, and
\eqref{prof:binary-transfer} give $e'\le Ce$ on each compact band.
Therefore $e=0$ through $T$.

The selection at the start of this paragraph licenses these
inequalities in integral form on every selected compact band.
The omitted time sets have vanishing measure, and the uniform
Lipschitz bound on \(e_n\) makes their contribution vanish.
Thus the limiting conclusion \(e=0\) applies to the selected
sequence of violating functions.

As $e(T)=0$, \eqref{prof:initial-entropy} gives $H_T\ge c'n-o(n)$.
During the terminal extension the bounded covariance input holds with
exponentially small exceptions even for an arbitrary input spin law.
It therefore applies directly to the $\nu$-field law.  The proof of
\eqref{prof:backward-entropy} on this fixed interval gives
$H_E\ge cH_T-o(n)\ge c''n$.  At the product endpoint it is enough to use
$\max_i|(h_E)_i|\le C_B\sqrt n$, with exception at most $e^{-Bn}$ for any
fixed $B$.  This follows from $\|R\|=O(1)$, $\|X\|=\sqrt n$, and Gaussian
tails.

Choose \(B\) to pay the restricted bound \(f^2\le100e^{2h}\).
The deterministic estimate in \eqref{prof:product-exception},
now with threshold \(C_B\sqrt n\), and the corresponding weighted
Gaussian tail control the correlated exceptional energy. The same product entropy estimate now
gives
\[
 H_E\le C\sqrt n\,\mathcal D_J(f)+o(1)
      \le CC_*n^{5/6}+o(1)=o(n),
\]
a contradiction.  This completes all three ranges and the proposition.
\end{proof}

\Needspace{12\baselineskip}
\subsection{Removing the bound on the test function}

\begin{theorem}[Uniform support profile]
\label{prof:support-profile}
\label{prof:thm-profile}
For every fixed $\omega>0$ there are $c_\omega>0$ and disorder events of
probability tending to one on which, simultaneously for every nonzero real
function $f$ with $p=\mu_J(f\ne0)\le1/2$,
\begin{equation}
 \mathcal D_J(f)
   \ge c_\omega n^{-2/3-\omega}\log(1/p)\,\mu_J(f^2).
 \label{prof:full-profile}
\end{equation}
In particular the spectral gap for rate-one-per-site heat-bath dynamics is
at least $c_\omega(\log2)n^{-2/3-\omega}$ on the same event.
\end{theorem}

\begin{proof}
The restricted profile holds for every admissible test on one disorder
event. We may therefore apply it to all dyadic pieces of an arbitrary
function on that same event. Pieces too small to meet the likelihood
bound will lose only a fixed fraction of the squared norm, and their
combined heat-bath energy is bounded by the original energy.

It suffices to consider $f\ge0$: replacing $f$ by $|f|$ preserves its
support and squared norm and cannot increase any conditional-variance
form.  For every integer $j$ put
\[
 g_j=\min\{(f-2^j)_+,2^j\},\qquad
 A_j=\mu(g_j^2),\qquad p_j=\mu(f>2^j).
\]
For a scalar $u\ge0$ the same truncations satisfy
\begin{equation}
 \frac{u^2}{4}\le
       \sum_{j\in\mathbb Z}\min\{(u-2^j)_+,2^j\}^2
       \le u^2.
 \label{prof:dyadic-norms}
\end{equation}
For $2^m\le u\le2^{m+1}$ the sum is
$4^m/3+(u-2^m)^2$, which verifies both bounds; $u=0$ is immediate.
Thus $S:=\sum_jA_j$ is finite and lies between $\mu(f^2)/4$ and
$\mu(f^2)$.

The truncation increments also obey, for $u\ge v\ge0$,
\[
 d_j:=\min\{(u-2^j)_+,2^j\}
      -\min\{(v-2^j)_+,2^j\}\ge0,
 \qquad \sum_jd_j=u-v.
\]
The sum identity follows by integrating the indicators of the disjoint
intervals $(2^j,2^{j+1})$.  Hence $\sum_jd_j^2\le(u-v)^2$.
Using the representation of conditional variance by two independent
conditional samples gives
\begin{equation}
 \sum_j\mathcal D_J(g_j)\le\mathcal D_J(f).
 \label{prof:dyadic-energies}
\end{equation}
All sums are justified by nonnegativity and monotone convergence.

Discard indices with $A_j<4^jp_j^2/100$, as well as $p_j=0$.
On $\{f>2^j\}$ the preceding piece $g_{j-1}$ is saturated, so
$A_{j-1}\ge4^{j-1}p_j$.  Since $p_j\le p\le1/2$, the discarded pieces
satisfy
\[
 \sum_{j\,\mathrm{discarded}}A_j
 \le\sum_j\frac{4^jp_j^2}{100}
 \le\frac1{50}\sum_jA_{j-1}=\frac S{50}.
\]
Every remaining normalized piece $g_j/\sqrt{A_j}$ has squared supremum at
most $4^j/A_j\le100/p_j^2$.  Proposition~\ref{prof:restricted-profile},
on its single event valid for all test functions, therefore gives
\[
 \begin{split}
 \mathcal D_J(f)
 &\ge c_\omega n^{-2/3-\omega}
       \sum_{j\,\mathrm{remaining}}A_j\log(1/p_j)\\
 &\ge\frac{49}{200}c_\omega n^{-2/3-\omega}
                 \log(1/p)\,\mu(f^2).
 \end{split}
\]
Renaming the constant proves \eqref{prof:full-profile}.

Finally, for an arbitrary real $u$ choose a $\mu$-median $m$, so that
$u_+=(u-m)_+$ and $u_-=(m-u)_+$ both have support mass at most $1/2$.
For each pair of values the squared increments of these two functions sum
to at most the squared increment of $u$.  Thus the support profile gives
\[
 \mathcal D_J(u)\ge\mathcal D_J(u_+)+\mathcal D_J(u_-)
 \ge c_\omega n^{-2/3-\omega}\log2\,\mu((u-m)^2)
 \ge c_\omega n^{-2/3-\omega}\log2\,
                 \operatorname{Var}_\mu(u).
\]
This is the asserted gap bound.
\end{proof}

The support profile and its Poincar\'e consequence complete the first
functional route. Section~\ref{sec:dynamics} converts this endpoint to
worst-start mixing in both time conventions.
